\documentclass[11pt,a4paper]{article}

\usepackage{iftex}
\ifXeTeX
  \usepackage{fontspec}
  \setmainfont{TeX Gyre Pagella}
  \setmonofont{DejaVu Sans Mono}[Scale=0.82]
\fi

\usepackage[a4paper, top=25mm, bottom=25mm, left=28mm, right=28mm]{geometry}
\usepackage{xcolor}
\definecolor{headingblue}{RGB}{31,56,100}
\definecolor{ruleblue}{RGB}{46,90,156}
\definecolor{codebg}{RGB}{245,245,245}

\usepackage{longtable,booktabs,array,calc,multirow}
\usepackage{microtype}
\usepackage{parskip}
\setlength{\parskip}{5pt}

\usepackage{fancyvrb}
\usepackage{fvextra}
\usepackage{framed}
\newenvironment{Shaded}{%
  \setlength{\FrameSep}{4pt}%
  \definecolor{shadecolor}{RGB}{245,245,245}%
  \begin{snugshade}\footnotesize
}{\end{snugshade}}
\DefineVerbatimEnvironment{Highlighting}{Verbatim}{%
  breaklines,breakanywhere,commandchars=\\\{\}}

\newcommand{\AlertTok}[1]{\textcolor[rgb]{0.94,0.16,0.16}{#1}}
\newcommand{\AnnotationTok}[1]{\textcolor[rgb]{0.56,0.35,0.01}{\textit{#1}}}
\newcommand{\AttributeTok}[1]{\textcolor[rgb]{0.13,0.29,0.53}{#1}}
\newcommand{\BaseNTok}[1]{\textcolor[rgb]{0.00,0.00,0.81}{#1}}
\newcommand{\BuiltInTok}[1]{#1}
\newcommand{\CharTok}[1]{\textcolor[rgb]{0.31,0.60,0.02}{#1}}
\newcommand{\CommentTok}[1]{\textcolor[rgb]{0.56,0.35,0.01}{\textit{#1}}}
\newcommand{\CommentVarTok}[1]{\textcolor[rgb]{0.56,0.35,0.01}{\textbf{\textit{#1}}}}
\newcommand{\ConstantTok}[1]{\textcolor[rgb]{0.56,0.35,0.01}{#1}}
\newcommand{\ControlFlowTok}[1]{\textcolor[rgb]{0.13,0.29,0.53}{\textbf{#1}}}
\newcommand{\DataTypeTok}[1]{\textcolor[rgb]{0.13,0.29,0.53}{#1}}
\newcommand{\DecValTok}[1]{\textcolor[rgb]{0.00,0.00,0.81}{#1}}
\newcommand{\DocumentationTok}[1]{\textcolor[rgb]{0.56,0.35,0.01}{\textbf{\textit{#1}}}}
\newcommand{\ErrorTok}[1]{\textcolor[rgb]{0.64,0.00,0.00}{\textbf{#1}}}
\newcommand{\ExtensionTok}[1]{#1}
\newcommand{\FloatTok}[1]{\textcolor[rgb]{0.00,0.00,0.81}{#1}}
\newcommand{\FunctionTok}[1]{\textcolor[rgb]{0.13,0.29,0.53}{#1}}
\newcommand{\ImportTok}[1]{#1}
\newcommand{\InformationTok}[1]{\textcolor[rgb]{0.56,0.35,0.01}{\textbf{\textit{#1}}}}
\newcommand{\KeywordTok}[1]{\textcolor[rgb]{0.13,0.29,0.53}{\textbf{#1}}}
\newcommand{\NormalTok}[1]{#1}
\newcommand{\OperatorTok}[1]{\textcolor[rgb]{0.81,0.36,0.00}{\textbf{#1}}}
\newcommand{\OtherTok}[1]{\textcolor[rgb]{0.56,0.35,0.01}{#1}}
\newcommand{\PreprocessorTok}[1]{\textcolor[rgb]{0.56,0.35,0.01}{\textit{#1}}}
\newcommand{\RegionMarkerTok}[1]{#1}
\newcommand{\SpecialCharTok}[1]{\textcolor[rgb]{0.81,0.36,0.00}{#1}}
\newcommand{\SpecialStringTok}[1]{\textcolor[rgb]{0.31,0.60,0.02}{#1}}
\newcommand{\StringTok}[1]{\textcolor[rgb]{0.31,0.60,0.02}{#1}}
\newcommand{\VariableTok}[1]{\textcolor[rgb]{0.00,0.00,0.00}{#1}}
\newcommand{\VerbatimStringTok}[1]{\textcolor[rgb]{0.31,0.60,0.02}{#1}}
\newcommand{\WarningTok}[1]{\textcolor[rgb]{0.56,0.35,0.01}{\textbf{\textit{#1}}}}
\providecommand{\tightlist}{\setlength{\itemsep}{2pt}\setlength{\parskip}{0pt}}

\usepackage{titlesec}
\titleformat{\section}{\large\bfseries\color{headingblue}}{\thesection.}{0.6em}{}[\vspace{-4pt}\color{ruleblue}\rule{\linewidth}{0.6pt}]
\titleformat{\subsection}{\normalsize\bfseries\color{headingblue}}{\thesubsection}{0.6em}{}
\titleformat{\subsubsection}{\normalsize\bfseries}{\thesubsubsection}{0.6em}{}
\titlespacing{\section}{0pt}{18pt}{6pt}
\titlespacing{\subsection}{0pt}{12pt}{4pt}
\titlespacing{\subsubsection}{0pt}{10pt}{3pt}
\setlength{\tabcolsep}{6pt}
\renewcommand{\arraystretch}{1.25}

\usepackage{hyperref}
\hypersetup{
  colorlinks=true,
  linkcolor=ruleblue,
  urlcolor=ruleblue,
  pdftitle={GNU Backgammon Android Port -- MASTER v0.9},
  pdfauthor={clavierhaus.at},
}

\usepackage{fancyhdr}
\pagestyle{fancy}
\fancyhf{}
\fancyhead[L]{\small\color{gray}GNU Backgammon Android Port --- MASTER v0.9}
\fancyhead[R]{\small\color{gray}clavierhaus.at}
\fancyfoot[C]{\small\color{gray}\thepage}
\renewcommand{\headrulewidth}{0.3pt}

\title{\textbf{GNU Backgammon Android Port}}
\author{clavierhaus.at}
\date{Juni 2026}

\begin{document}

\maketitle\thispagestyle{fancy}

\tableofcontents\newpage

\section{GNU Backgammon Android Port --- MASTER v0.9}\label{gnu-backgammon-android-port-master-v09}

\textbf{GNU Backgammon by clavierhaus.at}\\
clavierhaus.at · Vienna · Austria\\
June 2026

\begin{center}\rule{0.5\linewidth}{0.5pt}\end{center}

\subsection{1. Introduction \& Project
Objective}\label{introduction-project-objective}

The objective of this project is to produce a \textbf{true port} of GNU
Backgammon to the Android platform --- not a stripped-down evaluator,
but a complete port that mirrors 100\% of the functionality of the PC
version. This includes full game management, SGF import/export,
analysis, rollout, cube decisions, tutor mode, and the complete match
state machine.

The primary constraint is the monolithic nature of the original gnubg
desktop codebase, which is tightly coupled with GTK/GNOME UI components,
file I/O, global configuration state, and platform-specific threading
primitives. The approach is to compile the genuine gnubg engine source
--- every file that is portable --- and replace only what is inherently
GTK-specific with Android equivalents via a clean architectural
boundary: \texttt{android-app.c}.

\textbf{The key architectural insight:} gnubg's codebase separates
cleanly into three layers:

\begin{enumerate}
\def\labelenumi{\arabic{enumi}.}
\tightlist
\item
  \textbf{Mathematical engine} (\texttt{eval.c}, \texttt{rollout.c},
  \texttt{dice.c}, \texttt{bearoff.c} etc.) --- pure computation, zero
  UI dependencies, 100\% portable
\item
  \textbf{Game logic layer} (\texttt{play.c}, \texttt{analysis.c},
  \texttt{sgf.c}, \texttt{set.c}, \texttt{format.c} etc.) --- all GTK
  references are behind \texttt{\#if\ defined(USE\_GTK)} guards,
  portable when compiled without \texttt{USE\_GTK}
\item
  \textbf{GTK application shell} (\texttt{gnubg.c}, \texttt{progress.c},
  \texttt{gtk/*.c}) --- not portable, replaced by \texttt{android-app.c}
\end{enumerate}

Every stub that silently drops functionality is a hole in the port. The
goal is to minimise stubs to only what is genuinely GTK-specific
rendering code.

\subsubsection{1.1 Scope}\label{scope}

This document covers the complete build history, all architectural
decisions, every obstacle encountered and its resolution, and the exact
current state of the project as of Version 6. It is intended as a
working reference for any software engineer picking up or contributing
to this project.

\subsubsection{1.2 App Identity}\label{app-identity}

\begin{longtable}[]{@{}ll@{}}
\toprule\noalign{}
Property & Value \\
\midrule\noalign{}
\endhead
\bottomrule\noalign{}
\endlastfoot
App name & GNU Backgammon by clavierhaus.at \\
Android package & \texttt{com.clavierhaus.gnubg} \\
Minimum Android API & 31 (Android 12.0, Snow Cone) \\
Target ABI & \texttt{arm64-v8a} \\
Engine source & GNU Backgammon upstream (gnubg.org) \\
Repository & https://github.com/clavierhaus/gnubg-android \\
\end{longtable}

\begin{center}\rule{0.5\linewidth}{0.5pt}\end{center}

\subsection{2. Narrative of Progress}\label{narrative-of-progress}

\subsubsection{Phase 1: Monolithic Analysis \&
Prototype}\label{phase-1-monolithic-analysis-prototype}

Initial efforts focused on analysing the upstream GNU Backgammon source.
Attempts to build the full desktop environment proved unsustainable due
to deep dependencies on GTK, GLib, and global UI state variables
(\texttt{ms}, \texttt{ap}, \texttt{positions}). A ``Frankenstein'' build
approach --- attempting to stub these dependencies incrementally ---
established that the mathematical evaluation logic was sound in
isolation, but produced a fragile, non-portable build environment that
could not be taken further.

\subsubsection{Phase 2: Redesign --- Deterministic
Engineering}\label{phase-2-redesign-deterministic-engineering}

Phase 2 pivoted to a ``Clean-Room'' architecture. The attempt to port
the desktop environment was discarded entirely. The ``Core Engine'' was
isolated: a strict build hierarchy was established in which the upstream
source acts as a read-only reference, and the \texttt{engine-core/}
directory contains only the verified, essential mathematical components.

This phase established the orchestration scripts
(\texttt{pull\_dependencies.sh}, \texttt{patch\_engine.sh},
\texttt{build\_all.sh}) and the Patch Registry: surgical, idempotent
fixes for POSIX/Android compliance.

\subsubsection{Phase 3: The Pristine Headless Build
(Completed)}\label{phase-3-the-pristine-headless-build-completed}

Phase 3 achieved a fully hermetic host build of the mathematical core,
confirming that the engine compiles and links cleanly on Linux when
detached from the GTK/GNOME infrastructure.

When cross-compilation for Android NDK was attempted, the shim approach
collapsed. The root cause: gnubg uses GLib for threading primitives
(\texttt{GMutex}, \texttt{GCond}, \texttt{GThread}), dynamic data
structures (\texttt{GList}, \texttt{GArray}, \texttt{GHashTable}), and
character encoding conversion. Shimming these by hand created an
infinite chain of type collisions. This is not a sustainable path; it is
an attempt to reconstruct the Linux runtime inside the Android NDK.

\subsubsection{Phase 4: GLib Cross-Compilation for Android
(Completed)}\label{phase-4-glib-cross-compilation-for-android-completed}

The architectural decision was made to cross-compile a minimal but real
GLib 2.88.1 for \texttt{arm64-v8a} Android API 28, rather than shimming
GLib. This provides the full, tested GLib API that gnubg requires, with
no behavioral deviations from upstream.

The GLib source was obtained from the Fedora 44 SRPM
(\texttt{glib2-2.88.1-1.fc44}), ensuring Fedora's tested patches were
applied. The result is \texttt{libglib-2.0.so} verified as ELF 64-bit
ARM aarch64 for Android 28.

\subsubsection{Phase 5: JNI Bridge \& libgnubg-engine.so
(Completed)}\label{phase-5-jni-bridge-libgnubg-engine.so-completed}

Phase 5 constructed the JNI bridge layer (\texttt{native-lib.c},
\texttt{stubs.c}, \texttt{Engine.kt}) and CMake build pipeline producing
\texttt{libgnubg-engine.so}, verified as ELF 64-bit ARM aarch64 for
Android 28.

\subsubsection{Phase 6: NEON Acceleration \& Device Verification
(Completed)}\label{phase-6-neon-acceleration-device-verification-completed}

Phase 6 enabled ARM NEON SIMD acceleration and verified the engine
running correctly on a Pixel 8 Pro (aarch64, Android 16) via adb test
harness. Two critical bugs were discovered and fixed during device
testing; see §5.5 and §5.6.

\textbf{Verified results on device (opening position, 1-ply cubeless):}

\begin{longtable}[]{@{}lll@{}}
\toprule\noalign{}
Output & Value & Expected \\
\midrule\noalign{}
\endhead
\bottomrule\noalign{}
\endlastfoot
\texttt{ClassifyPosition} & 10 (CLASS\_CONTACT) & correct \\
Win probability & 0.5269 & \textasciitilde0.50 \\
Win gammon & 0.1478 & \textasciitilde0.12 \\
Win backgammon & 0.0085 & \textasciitilde0.01 \\
Lose gammon & 0.1285 & \textasciitilde0.12 \\
Lose backgammon & 0.0049 & \textasciitilde0.01 \\
\texttt{FindBestMove} 3-1 & 8/5 6/5 & correct (textbook) \\
\end{longtable}

\subsubsection{Phase 7: Cube Decisions \& Rollout
(Completed)}\label{phase-7-cube-decisions-rollout-completed}

Phase 7 added cube decision analysis and synchronous rollout to the JNI
API, both verified on device.

\textbf{Cube decision} (\texttt{Engine.cubeDecision()}): wraps
\texttt{GeneralCubeDecisionENoLocking} + \texttt{FindCubeDecision}.
Returns the full \texttt{aarOutput{[}2{]}{[}7{]}} equity matrix and the
\texttt{cubedecision} enum value. Verified: opening position returns
\texttt{NODOUBLE\_TAKE} (correct for a money game with a centred cube).

\textbf{Rollout} (\texttt{Engine.rollout()}): Phase 7 introduced a
synchronous single-threaded rollout via \texttt{gnubg\_rollout()}
calling \texttt{BasicCubefulRolloutNoLocking} directly. Phase 9
supersedes this with a full multi-threaded implementation; see Phase 9.
See §5.7 and §6.2.

\subsubsection{Phase 8: Full Game Logic Layer
(Completed)}\label{phase-8-full-game-logic-layer-completed}

Phase 8 is the pivotal phase --- the port moves from ``evaluator with
stubs'' to ``genuine gnubg port.'' The game logic layer
(\texttt{play.c}, \texttt{analysis.c}, \texttt{sgf.c} and all their
dependencies) is compiled into \texttt{libgnubg-engine.so}.

\textbf{The android-app.c architecture:} Rather than adding
\texttt{gnubg.c} (the GTK application shell, \textasciitilde8000 lines,
deeply GTK-dependent), a new file \texttt{jni-bridge/src/android-app.c}
was created as the Android equivalent. It contains: - Functions
extracted verbatim from \texttt{gnubg.c}: \texttt{InitBoard},
\texttt{NextToken}, \texttt{NextTokenGeneral}, \texttt{DisectPath},
\texttt{ParseNumber}, \texttt{ParseReal}, \texttt{ParsePosition},
\texttt{ParseKeyValue}, \texttt{ParsePlayer}, \texttt{SetToggle},
\texttt{CompareNames}, \texttt{GetMatchStateCubeInfo},
\texttt{find\_skills}, \texttt{swapGame}, \texttt{NameIsKey},
\texttt{AddKeyName}, \texttt{DeleteKeyName},
\texttt{CommandSwapPlayers}, \texttt{SmartSit}, \texttt{asyncEvalRoll},
\texttt{asyncMoveDecisionE}, \texttt{asyncFindMove},
\texttt{asyncCubeDecision}, \texttt{asyncAnalyzeMove},
\texttt{GetEvalChequer}, \texttt{GetEvalCube},
\texttt{GetEvalMoveFilter} - Android replacements for GTK functions:
\texttt{playSound→gnubg\_on\_sound\_event},
\texttt{ShowBoard→gnubg\_on\_board\_changed},
\texttt{UpdateSetting→gnubg\_on\_setting\_changed},
\texttt{CommandFirstGame/Move→gnubg\_on\_navigation\_event},
\texttt{GiveAdvice→gnubg\_on\_navigation\_event},
\texttt{RunAsyncProcess} (synchronous), \texttt{get\_input\_discard}
(returns TRUE), \texttt{confirmOverwrite} (returns TRUE) -
\texttt{RolloutProgressStart/Progress/End} replacing \texttt{progress.c}
(deeply GTK) - Sound infrastructure replacing \texttt{sound.c} - All
globals previously in \texttt{gnubg.c}

\textbf{Source files added to build:} \texttt{play.c},
\texttt{analysis.c}, \texttt{format.c}, \texttt{formatgs.c},
\texttt{drawboard.c}, \texttt{external.c}, \texttt{glib-ext.c},
\texttt{matchid.c}, \texttt{set.c}, \texttt{renderprefs.c},
\texttt{sgf.c}, \texttt{sgf\_y.c} (bison-generated), \texttt{sgf\_l.c}
(flex-generated)

\textbf{SGF parser generation:} \texttt{sgf\_y.y} and \texttt{sgf\_l.l}
are yacc/lex source files. They must be processed before building:

\begin{Shaded}
\begin{Highlighting}[]
\BuiltInTok{cd}\NormalTok{ engine{-}core}
\FunctionTok{bison} \AttributeTok{{-}d} \AttributeTok{{-}o}\NormalTok{ sgf\_y.c sgf\_y.y}
\FunctionTok{flex} \AttributeTok{{-}o}\NormalTok{ lex.yy.c sgf\_l.l }\KeywordTok{\&\&} \FunctionTok{mv}\NormalTok{ lex.yy.c sgf\_l.c}
\FunctionTok{sed} \AttributeTok{{-}i} \StringTok{\textquotesingle{}1s/\^{}/\#include \textless{}unistd.h\textgreater{}\textbackslash{}n/\textquotesingle{}}\NormalTok{ sgf\_l.c}
\end{Highlighting}
\end{Shaded}

\textbf{Verified:} All 6 existing test harness tests still pass on Pixel
8 Pro after adding the full game logic layer.

\subsubsection{Phase 9: Multi-threaded Rollout \& SGF JNI API
(Completed)}\label{phase-9-multi-threaded-rollout-sgf-jni-api-completed}

Phase 9 re-enables \texttt{USE\_MULTITHREAD} and implements
production-quality parallel rollout, then adds the SGF load/save JNI
entry points to complete the engine feature set.

\textbf{Multi-threaded rollout:} A persistent \texttt{GThreadPool} is
created at \texttt{gnubg\_init\_rollout()} with
\texttt{g\_get\_num\_processors()} threads. Each rollout trial is a
separate task dispatched via \texttt{g\_thread\_pool\_push()}. A
\texttt{GMutex}/\texttt{GCond} barrier blocks the calling thread until
all trials complete. Per-thread state is managed via two
\texttt{GPrivate} keys: \texttt{gnubg\_tls\_key} (owns
\texttt{ThreadLocalData} with \texttt{aMoves} and \texttt{pnnState}
buffers) and \texttt{gnubg\_rng\_key} (owns an isolated
\texttt{rngcontext} copy, preventing RNG data races). Result structs are
64-byte cache-line aligned to prevent false sharing on aarch64. The
scatter-gather merge computes mean and standard deviation after all
workers finish.

\textbf{SGF JNI entry points:} \texttt{Engine.loadSGF(path)} calls
\texttt{CommandLoadMatch()} which handles \texttt{FreeMatch},
\texttt{ClearMatch}, \texttt{RestoreGame}, and \texttt{UpdateSettings}
--- the complete match load sequence as implemented in \texttt{sgf.c}.
\texttt{Engine.saveSGF(path)} calls \texttt{CommandSaveMatch()} which
handles the \texttt{SaveGame} loop, file I/O,
\texttt{setDefaultFileName}, and \texttt{delete\_autosave}. Both entry
points are thread-safe via \texttt{gnubg\_lock}.

\textbf{Engine feature-complete:} After Phase 9,
\texttt{libgnubg-engine.so} provides 9 JNI entry points covering the
full gnubg evaluation and game management feature set. See §6.2 for the
complete API reference.

\begin{center}\rule{0.5\linewidth}{0.5pt}\end{center}

\subsection{3. Architecture}\label{architecture}

\subsubsection{3.1 Directory Structure}\label{directory-structure}

All project artefacts reside under
\texttt{/home/erweitert/gnubg-android/}:

\begin{verbatim}
/home/erweitert/gnubg-android/
  android-sdk/               <- Android SDK & NDK (NDK 27.0.11718014)
  engine-core/               <- gnubg source (config.h is Android-patched)
    lib/                     <- gnubg math library (neuralnet, SFMT, cache...)
    config.h                 <- PATCHED: host-generated flags removed
    lib/config.h             <- Wrapper: #include "../config.h"
  upstream-source/gnubg/     <- Unmodified upstream reference
  glib-2.88.1/               <- GLib source (Fedora SRPM, patched)
  glib-android-build/        <- Meson build directory (out-of-tree)
  jni-bridge/                <- JNI bridge layer
    CMakeLists.txt           <- NDK CMake build
    config.h                 <- Shadow config.h for engine-core/*.c
    external/glib/           <- Installed GLib (headers + .so files)
    src/native-lib.c         <- JNI entry points
    src/stubs.c              <- UI/threading layer stubs + TLD init
    src/com/clavierhaus/gnubg/Engine.kt  <- Kotlin JNI declarations
    build/libgnubg-engine.so <- Final deliverable
  test-harness/              <- Android adb test harness
    CMakeLists.txt           <- Builds standalone Android executable
    src/main.c               <- Test entry point
    src/stubs.c              <- Same stubs as jni-bridge
  android-arm64.cross        <- Meson cross-file for GLib build
  build_glib_android.sh      <- GLib cross-build script
  doc/                       <- Documentation
    Makefile                 <- make / make pdf / make tex
    gnubg.tex                <- XeLaTeX template
\end{verbatim}

\subsubsection{3.2 Build Toolchain}\label{build-toolchain}

\begin{longtable}[]{@{}ll@{}}
\toprule\noalign{}
Component & Version / Detail \\
\midrule\noalign{}
\endhead
\bottomrule\noalign{}
\endlastfoot
Host OS & Fedora 44, KDE Plasma 6/Wayland \\
Compiler (host) & GCC 16.1.1 (Red Hat) \\
Compiler (Android) & Clang 18.0.1 (NDK r27-beta1, build 11718014) \\
NDK & 27.0.11718014 at \texttt{android-sdk/ndk/27.0.11718014/} \\
Meson & 1.11.1 (Fedora package) \\
Ninja & 1.13.2 (Fedora package) \\
CMake & 4.3.0 (Fedora package) \\
GLib & 2.88.1 (Fedora SRPM \texttt{glib2-2.88.1-1.fc44}) \\
Target ABI & \texttt{arm64-v8a} (\texttt{aarch64-linux-android28}) \\
Android API level & 28 (Android 9.0 Pie) \\
Test device & Pixel 8 Pro (aarch64, Android 16) \\
\end{longtable}

\begin{center}\rule{0.5\linewidth}{0.5pt}\end{center}

\subsection{4. GLib Cross-Compilation for
Android}\label{glib-cross-compilation-for-android}

\subsubsection{4.1 Rationale}\label{rationale}

gnubg's engine uses GLib for: threading primitives (\texttt{GMutex},
\texttt{GCond}, \texttt{GThread}); dynamic data structures
(\texttt{GList}, \texttt{GArray}, \texttt{GHashTable}); character
encoding conversion (iconv wrapper); memory allocation wrappers;
assertion macros; and logging. The engine cannot function without these.
Shimming them by hand is not viable: each shim requires further GLib
internals, leading to infinite type collision regression.

Cross-compiling real GLib for Android provides the full, tested
implementation that gnubg was written against, with no behavioral
deviation. The investment is a one-time build; the resulting
\texttt{.so} files are reused for all future engine builds.

\subsubsection{4.2 Meson Cross-File
(android-arm64.cross)}\label{meson-cross-file-android-arm64.cross}

The cross-file specifies the NDK toolchain binaries, target machine
description, and platform properties that Meson cannot auto-detect in a
cross build. Critical design points:

\begin{itemize}
\tightlist
\item
  \textbf{Do NOT define \texttt{\_\_ANDROID\_API\_\_} in
  \texttt{c\_args}:} The compiler binary
  \texttt{aarch64-linux-android28-clang} already encodes API level 28.
  NDK clang 18 defines it internally; a redefinition causes
  \texttt{-Werror} failures in Meson's size detection probes.
\item
  \textbf{Explicit \texttt{sizeof\_*} properties:} Provided to bypass
  Meson's cross-detection probes, which fail because they attempt to run
  target binaries on the host.
\item
  \textbf{\texttt{pkg-config} disabled:} Set to \texttt{false} to
  prevent Meson from accidentally finding host-side \texttt{.pc} files.
\end{itemize}

\subsubsection{4.3 Key Build Obstacles \&
Resolutions}\label{key-build-obstacles-resolutions}

\begin{longtable}[]{@{}
  >{\raggedright\arraybackslash}p{(\linewidth - 2\tabcolsep) * \real{0.5000}}
  >{\raggedright\arraybackslash}p{(\linewidth - 2\tabcolsep) * \real{0.5000}}@{}}
\toprule\noalign{}
\begin{minipage}[b]{\linewidth}\raggedright
Obstacle
\end{minipage} & \begin{minipage}[b]{\linewidth}\raggedright
Resolution
\end{minipage} \\
\midrule\noalign{}
\endhead
\bottomrule\noalign{}
\endlastfoot
Unknown option \texttt{iconv} & GLib 2.88 removed \texttt{-Diconv} as a
Meson option. iconv detection is now via
\texttt{dependency(\textquotesingle{}iconv\textquotesingle{})}
internally. \\
\texttt{size\_t} detection failure & Caused by
\texttt{-D\_\_ANDROID\_API\_\_=26} in \texttt{c\_args} conflicting with
the compiler's built-in definition. Removed the redundant define. \\
iconv not found (API 26) & Bionic declares iconv under
\texttt{\#if\ \_\_ANDROID\_API\_\_\ \textgreater{}=\ 28}. Bumped target
from API 26 to API 28, which is the correct minimum for this project. \\
\texttt{dependency(\textquotesingle{}iconv\textquotesingle{})} fails on
Android & GLib's Meson uses \texttt{clang++} to probe
\texttt{iconv\_open()}, but Bionic's \texttt{iconv.h} lacks
\texttt{extern\ "C"} linkage for C++ probing. Patched GLib's
\texttt{meson.build} to use \texttt{declare\_dependency()} on Android
since iconv is in Bionic libc with no separate library. \\
\texttt{libintl.h} not found & GLib's \texttt{gi18n.h} requires
\texttt{libintl.h}. The GLib build's \texttt{proxy-libintl} subproject
provides both the header and \texttt{libintl.so}. Added
\texttt{external/glib/include} to the CMake include path. \\
\end{longtable}

\subsubsection{4.4 Installed Artefacts}\label{installed-artefacts}

The GLib build installs to \texttt{jni-bridge/external/glib/} and
provides:

\begin{itemize}
\tightlist
\item
  \texttt{lib/libglib-2.0.so} --- core GLib (the primary dependency)
\item
  \texttt{lib/libintl.so} --- proxy-libintl stub (required by
  \texttt{gi18n.h} / gettext macros)
\item
  \texttt{lib/libffi.so} --- libffi (GLib closure support)
\item
  \texttt{lib/libgio-2.0.so}, \texttt{libgobject-2.0.so},
  \texttt{libgmodule-2.0.so}, \texttt{libgthread-2.0.so} --- GLib
  sublibraries
\item
  \texttt{include/glib-2.0/} --- GLib headers
\item
  \texttt{include/libintl.h} --- proxy-libintl header
\item
  \texttt{lib/glib-2.0/include/glibconfig.h} --- platform-specific GLib
  configuration
\end{itemize}

\begin{center}\rule{0.5\linewidth}{0.5pt}\end{center}

\subsection{5. JNI Bridge Construction}\label{jni-bridge-construction}

\subsubsection{5.1 Source File Inventory}\label{source-file-inventory}

\begin{longtable}[]{@{}
  >{\raggedright\arraybackslash}p{(\linewidth - 2\tabcolsep) * \real{0.5000}}
  >{\raggedright\arraybackslash}p{(\linewidth - 2\tabcolsep) * \real{0.5000}}@{}}
\toprule\noalign{}
\begin{minipage}[b]{\linewidth}\raggedright
File
\end{minipage} & \begin{minipage}[b]{\linewidth}\raggedright
Purpose
\end{minipage} \\
\midrule\noalign{}
\endhead
\bottomrule\noalign{}
\endlastfoot
\texttt{jni-bridge/CMakeLists.txt} & NDK CMake build definition. Lists
all compiled source files, include paths, compile definitions, and
linked libraries. \\
\texttt{jni-bridge/config.h} & Shadow \texttt{config.h} for
\texttt{engine-core/*.c} files. See §5.3. \\
\texttt{jni-bridge/format.h} & Empty stub header --- \texttt{rollout.c}
includes it but uses no symbols. Prevents GTK chain. \\
\texttt{jni-bridge/matchid.h} & Empty stub header --- same rationale as
\texttt{format.h}. \\
\texttt{jni-bridge/analysis.h} & Minimal stub --- \texttt{sgf.c}
includes it; \texttt{RestoreDoubleAnalysis}/\texttt{RestoreMoveAnalysis}
are defined within \texttt{sgf.c} itself. \\
\texttt{jni-bridge/drawboard.h} & Stub defining
\texttt{FORMATEDMOVESIZE\ 29}. Prevents GTK rendering chain. \\
\texttt{jni-bridge/render.h} & Minimal stub for \texttt{renderprefs.h}
include chain. \\
\texttt{jni-bridge/renderprefs.h} & Minimal stub --- \texttt{play.c}
includes it but uses no rendering symbols. \\
\texttt{jni-bridge/sound.h} & Not needed --- real
\texttt{engine-core/sound.h} exists and is used. \\
\texttt{engine-core/config.h} & Host-generated autotools config, patched
to remove flags invalid on Android. \\
\texttt{engine-core/lib/config.h} & Single-line wrapper:
\texttt{\#include\ "../config.h"}. \\
\texttt{engine-core/lib/neuralnetsse.c} & Modified: VLA replaced with
\texttt{posix\_memalign}. See §5.6. \\
\texttt{jni-bridge/src/stubs.c} & Remaining GTK/UI stubs, TLD init,
rollout infrastructure. See §5.4, §5.5, §5.7. \\
\texttt{jni-bridge/src/android-app.c} & \textbf{Android application
shell} replacing \texttt{gnubg.c}. See §5.8. \\
\texttt{jni-bridge/src/native-lib.c} & JNI entry points for all Engine
methods. \\
\texttt{jni-bridge/src/com/clavierhaus/gnubg/Engine.kt} & Kotlin object
declaring all external (JNI) functions. \\
\end{longtable}

\subsubsection{5.2 Engine Source Files in
Build}\label{engine-source-files-in-build}

The following \texttt{engine-core} source files are compiled into
\texttt{libgnubg-engine.so}:

\begin{longtable}[]{@{}
  >{\raggedright\arraybackslash}p{(\linewidth - 2\tabcolsep) * \real{0.5000}}
  >{\raggedright\arraybackslash}p{(\linewidth - 2\tabcolsep) * \real{0.5000}}@{}}
\toprule\noalign{}
\begin{minipage}[b]{\linewidth}\raggedright
File
\end{minipage} & \begin{minipage}[b]{\linewidth}\raggedright
Role
\end{minipage} \\
\midrule\noalign{}
\endhead
\bottomrule\noalign{}
\endlastfoot
\texttt{eval.c} & Evaluation engine core, neural network inference, move
generation \\
\texttt{dice.c} & Dice RNG (SFMT/Mersenne Twister; GMP/BBS and
random.org disabled) \\
\texttt{bearoff.c}, \texttt{bearoffgammon.c} & Bearoff database \\
\texttt{positionid.c} & Position ID encoding/decoding \\
\texttt{matchequity.c}, \texttt{mec.c} & Match equity tables \\
\texttt{util.c}, \texttt{file.c}, \texttt{boardpos.c} & Utility
functions \\
\texttt{multithread.c} & Threading infrastructure (single-threaded on
Android) \\
\texttt{rollout.c} & Monte Carlo rollout engine \\
\texttt{sgf.c} & SGF file handling \\
\texttt{sgf\_y.c} & SGF parser (bison-generated from
\texttt{sgf\_y.y}) \\
\texttt{sgf\_l.c} & SGF lexer (flex-generated from \texttt{sgf\_l.l}) \\
\texttt{play.c} & Game state machine, match management, move records \\
\texttt{analysis.c} & Position analysis, skill classification \\
\texttt{format.c} & Move and equity formatting \\
\texttt{formatgs.c} & Game statistics formatting \\
\texttt{drawboard.c} & Board ASCII rendering and move formatting
(\texttt{FormatMove}) \\
\texttt{external.c} & External player protocol (POSIX sockets) \\
\texttt{glib-ext.c} & GLib extension utilities (GValue/GObject
helpers) \\
\texttt{matchid.c} & Match ID encoding/decoding \\
\texttt{set.c} & Settings commands and configuration \\
\texttt{renderprefs.c} & Render preferences (GTK sections guarded) \\
\texttt{lib/output.c} & Output formatting \\
\texttt{lib/SFMT.c} & SIMD-oriented Fast Mersenne Twister \\
\texttt{lib/neuralnet.c} & Neural network evaluation \\
\texttt{lib/neuralnetsse.c} & NEON neural net implementation
(patched) \\
\texttt{lib/cache.c}, \texttt{lib/md5.c}, \texttt{lib/isaac.c},
\texttt{lib/list.c} & Supporting data structures \\
\texttt{lib/inputs.c} & Neural net input feature computation \\
\end{longtable}

\subsubsection{5.3 The config.h Shadow
Mechanism}\label{the-config.h-shadow-mechanism}

This is the most architecturally significant decision in the JNI bridge
build.

The upstream \texttt{engine-core/config.h} is generated by GNU autotools
\texttt{./configure} on the Fedora host. It contains \texttt{\#define}
flags for features present on Fedora that do not exist on Android:

\begin{Shaded}
\begin{Highlighting}[]
\PreprocessorTok{\#define HAVE\_LIBGMP }\DecValTok{1}\PreprocessorTok{            }\CommentTok{// GMP (dice.c BBS RNG)}
\PreprocessorTok{\#define USE\_SIMD\_INSTRUCTIONS }\DecValTok{1}\PreprocessorTok{  }\CommentTok{// x86 SSE/AVX — invalid on aarch64}
\PreprocessorTok{\#define HAVE\_SSE }\DecValTok{1}
\PreprocessorTok{\#define USE\_SSE2 }\DecValTok{1}
\PreprocessorTok{\#define USE\_AVX }\DecValTok{1}
\PreprocessorTok{\#define LIBCURL\_PROTOCOL\_HTTPS }\DecValTok{1}\PreprocessorTok{ }\CommentTok{// randomorg.c}
\PreprocessorTok{\#define HAVE\_LIBCURL }\DecValTok{1}
\PreprocessorTok{\#define USE\_MULTITHREAD }\DecValTok{1}\PreprocessorTok{        }\CommentTok{// GLib{-}based thread pool}
\end{Highlighting}
\end{Shaded}

These flags cannot be overridden via \texttt{-U} on the compiler command
line because the C standard specifies that \texttt{\#include\ "file.h"}
(quoted include) searches the source file's own directory before any
\texttt{-I} paths. This means \texttt{eval.c}'s
\texttt{\#include\ "config.h"} resolves to \texttt{engine-core/config.h}
regardless of any \texttt{-I} flags pointing elsewhere.

\textbf{The solution --- the shadow file mechanism:} a file named
\texttt{config.h} is placed in each directory from which source files
perform a quoted \texttt{config.h} include. Each shadow file includes
the real \texttt{config.h} via a relative path, undefines the invalid
flags, and adds Android-specific defines.

Two shadow files are required:

\begin{itemize}
\tightlist
\item
  \texttt{jni-bridge/config.h} --- intercepts includes from
  \texttt{engine-core/*.c}
\item
  \texttt{engine-core/lib/config.h} --- intercepts includes from
  \texttt{engine-core/lib/*.c}
\end{itemize}

The current shadow config adds NEON defines after stripping x86 SIMD:

\begin{Shaded}
\begin{Highlighting}[]
\PreprocessorTok{\#include }\ImportTok{"../engine{-}core/config.h"}
\PreprocessorTok{\#undef HAVE\_LIBGMP}
\PreprocessorTok{\#undef LIBCURL\_PROTOCOL\_HTTPS}
\PreprocessorTok{\#undef HAVE\_LIBCURL}
\PreprocessorTok{\#undef USE\_SIMD\_INSTRUCTIONS   }\CommentTok{// strip x86 SIMD first}
\PreprocessorTok{\#undef HAVE\_SSE}
\PreprocessorTok{\#undef USE\_SSE2}
\PreprocessorTok{\#undef USE\_AVX}
\CommentTok{// ARM NEON — mandatory on aarch64, replaces x86 SIMD:}
\PreprocessorTok{\#define USE\_SIMD\_INSTRUCTIONS }\DecValTok{1}
\PreprocessorTok{\#define HAVE\_NEON }\DecValTok{1}
\PreprocessorTok{\#define USE\_NEON }\DecValTok{1}
\end{Highlighting}
\end{Shaded}

The real \texttt{engine-core/config.h} is also regenerated without the
invalid flags using \texttt{grep\ -v}, making the shadow undefines
redundant but retained as defense-in-depth. \textbf{This approach
requires no modification to any upstream source file} (except
\texttt{neuralnetsse.c}; see §5.6).

\subsubsection{5.4 stubs.c Design}\label{stubs.c-design}

\paragraph{5.4.1 Global Variable Storage}\label{global-variable-storage}

Several gnubg globals are declared \texttt{extern} in headers but their
storage is defined in desktop-layer source files not in the Android
build. \texttt{stubs.c} provides the storage with correct types:

\begin{itemize}
\tightlist
\item
  \texttt{matchstate\ ms} --- the global match state struct
\item
  \texttt{player\ ap{[}2{]}} --- the two player structs
\item
  \texttt{int\ positions{[}2{]}{[}30{]}{[}3{]}} --- board position
  geometry array (from \texttt{boardpos.h})
\item
  \texttt{ThreadData\ td} --- the thread pool state struct
  (zero-initialised at declaration; populated by
  \texttt{gnubg\_init\_tld()})
\item
  \texttt{const\ char\ *szHomeDirectory},
  \texttt{char\ *szCurrentFileName} --- path globals
\end{itemize}

\paragraph{5.4.2 Function Stubs}\label{function-stubs}

Functions belonging to the GTK/UI/desktop layer are stubbed with
signatures matching the header declarations exactly:

\begin{itemize}
\tightlist
\item
  \textbf{Threading:} \texttt{Mutex\_Lock}, \texttt{Mutex\_Release},
  \texttt{ResetManualEvent}, \texttt{SetManualEvent},
  \texttt{WaitForManualEvent}, \texttt{InitManualEvent},
  \texttt{FreeManualEvent}, \texttt{InitMutex}, \texttt{CloseThread},
  \texttt{TLSSetValue}, \texttt{MT\_CreateThreadLocalData}
\item
  \textbf{Locking wrappers (EXP\_LOCK\_FUN):}
  \texttt{EvaluatePositionWithLocking},
  \texttt{FindBestMoveWithLocking},
  \texttt{FindnSaveBestMovesWithLocking},
  \texttt{GeneralCubeDecisionEWithLocking},
  \texttt{GeneralEvaluationEWithLocking}, \texttt{ScoreMoveWithLocking},
  \texttt{BasicCubefulRolloutWithLocking} --- each delegates to the
  corresponding \texttt{NoLocking} variant (single-threaded evaluation)
\item
  \textbf{UI/progress:} \texttt{ProcessEvents}, \texttt{ProgressValue},
  \texttt{ProgressStart}, \texttt{ProgressEnd},
  \texttt{ProgressValueAdd}
\item
  \textbf{Callbacks:} \texttt{LogCube}, \texttt{SetRNG},
  \texttt{ChangeGame}, \texttt{FormatMove},
  \texttt{get\_current\_moverecord}
\item
  \textbf{Network dice:} \texttt{RandomorgDice}, \texttt{NetworkDice}
  (returns -1; requires libcurl)
\item
  \textbf{Match state:} \texttt{save\_autosave}, \texttt{msBoard}
\end{itemize}

\paragraph{5.4.3 Disabled Features}\label{disabled-features}

\begin{longtable}[]{@{}
  >{\raggedright\arraybackslash}p{(\linewidth - 4\tabcolsep) * \real{0.3333}}
  >{\raggedright\arraybackslash}p{(\linewidth - 4\tabcolsep) * \real{0.3333}}
  >{\raggedright\arraybackslash}p{(\linewidth - 4\tabcolsep) * \real{0.3333}}@{}}
\toprule\noalign{}
\begin{minipage}[b]{\linewidth}\raggedright
Feature
\end{minipage} & \begin{minipage}[b]{\linewidth}\raggedright
Reason disabled
\end{minipage} & \begin{minipage}[b]{\linewidth}\raggedright
Re-enablement
\end{minipage} \\
\midrule\noalign{}
\endhead
\bottomrule\noalign{}
\endlastfoot
GMP/BBS RNG (\texttt{dice.c}) & Requires \texttt{libgmp}, unavailable on
Android & Add libgmp cross-build; restore \texttt{HAVE\_LIBGMP} in
\texttt{config.h} \\
x86 SIMD & \texttt{HAVE\_SSE}, \texttt{USE\_SSE2}, \texttt{USE\_AVX}
invalid on aarch64 & N/A --- replaced by NEON \\
GNU Multithread pool & Previously disabled; \texttt{USE\_MULTITHREAD}
now active & Re-enabled in Phase 9 via
\texttt{GThreadPool}/\texttt{GPrivate} \\
libcurl / random.org & \texttt{LIBCURL\_PROTOCOL\_HTTPS} /
\texttt{HAVE\_LIBCURL} removed; \texttt{randomorg.c} excluded &
Cross-compile libcurl for Android; restore flags and add
\texttt{randomorg.c} to build \\
\end{longtable}

\subsubsection{5.5 Thread-Local Data Initialisation
(Critical)}\label{thread-local-data-initialisation-critical}

\textbf{Background:} When \texttt{USE\_MULTITHREAD} is \textbf{enabled}
(Phase 9+), gnubg's move generation and scoring macros expand as
follows:

\begin{Shaded}
\begin{Highlighting}[]
\PreprocessorTok{\#define MT\_Get\_aMoves}\OperatorTok{()}\PreprocessorTok{   }\OperatorTok{((}\PreprocessorTok{ThreadLocalData }\OperatorTok{*)}\PreprocessorTok{TLSGet}\OperatorTok{(}\PreprocessorTok{td}\OperatorTok{.}\PreprocessorTok{tlsItem}\OperatorTok{)){-}\textgreater{}}\PreprocessorTok{aMoves}
\PreprocessorTok{\#define MT\_Get\_nnState}\OperatorTok{()}\PreprocessorTok{  }\OperatorTok{((}\PreprocessorTok{ThreadLocalData }\OperatorTok{*)}\PreprocessorTok{TLSGet}\OperatorTok{(}\PreprocessorTok{td}\OperatorTok{.}\PreprocessorTok{tlsItem}\OperatorTok{)){-}\textgreater{}}\PreprocessorTok{pnnState}
\end{Highlighting}
\end{Shaded}

Both call \texttt{TLSGet()} which is implemented via \texttt{GPrivate}
--- each thread gets its own \texttt{ThreadLocalData} allocated lazily
on first access. This replaces the previous single-threaded approach of
a single static \texttt{gnubg\_tld} struct.

\begin{Shaded}
\begin{Highlighting}[]
\PreprocessorTok{\#define MT\_Get\_aMoves}\OperatorTok{()}\PreprocessorTok{   td}\OperatorTok{.}\PreprocessorTok{tld}\OperatorTok{{-}\textgreater{}}\PreprocessorTok{aMoves}
\PreprocessorTok{\#define MT\_Get\_nnState}\OperatorTok{()}\PreprocessorTok{  td}\OperatorTok{.}\PreprocessorTok{tld}\OperatorTok{{-}\textgreater{}}\PreprocessorTok{pnnState}
\end{Highlighting}
\end{Shaded}

Both dereference \texttt{td.tld}, a \texttt{ThreadLocalData\ *} inside
the global \texttt{ThreadData\ td}. At program start, \texttt{td} is
zero-initialised, making \texttt{td.tld\ =\ NULL}. The first call to
1-ply evaluation triggers \texttt{GenerateMoves} →
\texttt{MT\_Get\_aMoves()} → null dereference at offset 0x8. Similarly,
\texttt{ScoreMoves} calls \texttt{MT\_Get\_nnState()} →
\texttt{td.tld-\textgreater{}pnnState}, also null, causing a write
through null at offset 0x40 into the NNState array.

\textbf{The fix --- \texttt{gnubg\_init\_tld()}:} A function in
\texttt{stubs.c} that must be called once after
\texttt{EvalInitialise()} and before any evaluation:

\begin{Shaded}
\begin{Highlighting}[]
\DataTypeTok{static}\NormalTok{ ThreadLocalData gnubg\_tld}\OperatorTok{;}
\DataTypeTok{static}\NormalTok{ NNState         gnubg\_nn\_states}\OperatorTok{[}\DecValTok{3}\OperatorTok{];}
\DataTypeTok{static}\NormalTok{ move            gnubg\_moves}\OperatorTok{[}\NormalTok{MAX\_MOVES}\OperatorTok{];}

\DataTypeTok{void}\NormalTok{ gnubg\_init\_tld}\OperatorTok{(}\DataTypeTok{void}\OperatorTok{)} \OperatorTok{\{}
    \DataTypeTok{unsigned} \DataTypeTok{int}\NormalTok{ i}\OperatorTok{;}

\NormalTok{    gnubg\_tld}\OperatorTok{.}\NormalTok{aMoves    }\OperatorTok{=}\NormalTok{ gnubg\_moves}\OperatorTok{;}
\NormalTok{    gnubg\_tld}\OperatorTok{.}\NormalTok{pnnState  }\OperatorTok{=}\NormalTok{ gnubg\_nn\_states}\OperatorTok{;}
\NormalTok{    td}\OperatorTok{.}\NormalTok{tld              }\OperatorTok{=} \OperatorTok{\&}\NormalTok{gnubg\_tld}\OperatorTok{;}

    \CommentTok{/* Allocate savedBase (cHidden floats) and savedIBase (cInput floats)}
\CommentTok{     * for each NNState entry. Sizes are read from nnContact after weight load.}
\CommentTok{     * Three entries cover CLASS\_RACE, CLASS\_CRASHED, CLASS\_CONTACT offsets. */}
    \ControlFlowTok{for} \OperatorTok{(}\NormalTok{i }\OperatorTok{=} \DecValTok{0}\OperatorTok{;}\NormalTok{ i }\OperatorTok{\textless{}} \DecValTok{3}\OperatorTok{;}\NormalTok{ i}\OperatorTok{++)} \OperatorTok{\{}
\NormalTok{        gnubg\_nn\_states}\OperatorTok{[}\NormalTok{i}\OperatorTok{].}\NormalTok{state     }\OperatorTok{=}\NormalTok{ NNSTATE\_NONE}\OperatorTok{;}
\NormalTok{        gnubg\_nn\_states}\OperatorTok{[}\NormalTok{i}\OperatorTok{].}\NormalTok{savedBase  }\OperatorTok{=}\NormalTok{ g\_malloc}\OperatorTok{(}\NormalTok{nnContact}\OperatorTok{.}\NormalTok{cHidden }\OperatorTok{*} \KeywordTok{sizeof}\OperatorTok{(}\DataTypeTok{float}\OperatorTok{));}
\NormalTok{        gnubg\_nn\_states}\OperatorTok{[}\NormalTok{i}\OperatorTok{].}\NormalTok{savedIBase }\OperatorTok{=}\NormalTok{ g\_malloc}\OperatorTok{(}\NormalTok{nnContact}\OperatorTok{.}\NormalTok{cInput  }\OperatorTok{*} \KeywordTok{sizeof}\OperatorTok{(}\DataTypeTok{float}\OperatorTok{));}
    \OperatorTok{\}}
\OperatorTok{\}}
\end{Highlighting}
\end{Shaded}

\textbf{Why 3 NNState entries:} \texttt{ScoreMoves} accesses
\texttt{nnStates{[}0..2{]}} directly, and \texttt{EvaluatePositionFull}
indexes by \texttt{pc\ -\ CLASS\_RACE} where \texttt{pc} ranges from
\texttt{CLASS\_RACE} (8) to \texttt{CLASS\_CONTACT} (10), giving offsets
0, 1, 2.

\textbf{Call site in native-lib.c:} \texttt{gnubg\_init\_tld()} must be
called inside \texttt{Java\_com\_clavierhaus\_gnubg\_Engine\_initialise}
after \texttt{EvalInitialise()} returns, while \texttt{gnubg\_lock} is
held.

\textbf{Call site in test harness:} Called immediately after
\texttt{EvalInitialise()} in \texttt{main()}.

\subsubsection{5.6 NeuralNetEvaluateSSE VLA Alignment
Patch}\label{neuralnetevaluatesse-vla-alignment-patch}

\textbf{Background:} \texttt{NeuralNetEvaluateSSE} in
\texttt{neuralnetsse.c} declares its hidden-layer activation buffer as a
Variable Length Array with an alignment attribute:

\begin{Shaded}
\begin{Highlighting}[]
\NormalTok{SSE\_ALIGN}\OperatorTok{(}\DataTypeTok{float}\NormalTok{ ar}\OperatorTok{[}\NormalTok{pnn}\OperatorTok{{-}\textgreater{}}\NormalTok{cHidden}\OperatorTok{]);}
\CommentTok{// expands to: float ar[pnn{-}\textgreater{}cHidden] \_\_attribute\_\_((aligned(16)));}
\end{Highlighting}
\end{Shaded}

On aarch64 with Clang 18, the \texttt{\_\_attribute\_\_((aligned(16)))}
on a VLA is not guaranteed to be honoured. The NEON intrinsics
(\texttt{vld1q\_f32}, \texttt{vst1q\_f32}) then operate on a potentially
misaligned buffer, causing a SIGSEGV.

\textbf{The fix:} Replace the VLA with a heap-allocated aligned buffer:

\begin{Shaded}
\begin{Highlighting}[]
\DataTypeTok{float} \OperatorTok{*}\NormalTok{ar }\OperatorTok{=}\NormalTok{ NULL}\OperatorTok{;}
\ControlFlowTok{if} \OperatorTok{(}\NormalTok{posix\_memalign}\OperatorTok{((}\DataTypeTok{void} \OperatorTok{**)\&}\NormalTok{ar}\OperatorTok{,}\NormalTok{ ALIGN\_SIZE}\OperatorTok{,}\NormalTok{ pnn}\OperatorTok{{-}\textgreater{}}\NormalTok{cHidden }\OperatorTok{*} \KeywordTok{sizeof}\OperatorTok{(}\DataTypeTok{float}\OperatorTok{))} \OperatorTok{!=} \DecValTok{0}\OperatorTok{)}
    \ControlFlowTok{return} \OperatorTok{{-}}\DecValTok{1}\OperatorTok{;}
\NormalTok{EvaluateSSE}\OperatorTok{(}\NormalTok{pnn}\OperatorTok{,}\NormalTok{ arInput}\OperatorTok{,}\NormalTok{ ar}\OperatorTok{,}\NormalTok{ arOutput}\OperatorTok{);}
\NormalTok{free}\OperatorTok{(}\NormalTok{ar}\OperatorTok{);}
\ControlFlowTok{return} \DecValTok{0}\OperatorTok{;}
\end{Highlighting}
\end{Shaded}

\texttt{posix\_memalign} guarantees \texttt{ALIGN\_SIZE} (16 bytes for
NEON) alignment. This is the only modification made to a file in
\texttt{engine-core/lib/} beyond the shadow \texttt{config.h}. It is
documented in \texttt{PROVENANCE.md}.

\subsubsection{5.7 Rollout Infrastructure in
stubs.c}\label{rollout-infrastructure-in-stubs.c}

\paragraph{5.7.1 Evolution}\label{evolution}

Phase 7 introduced a synchronous single-threaded rollout bypassing
\texttt{MT\_WaitForTasks}. Phase 9 replaced this with a full
multi-threaded implementation using \texttt{GThreadPool}. Both share the
same external API (\texttt{gnubg\_rollout()}) and the same
\texttt{gnubg\_init\_rollout()} initialisation call.

\paragraph{5.7.2 Multi-threaded Architecture (Phase
9)}\label{multi-threaded-architecture-phase-9}

\texttt{gnubg\_rollout()} orchestrates a parallel rollout using a
persistent \texttt{GThreadPool}:

\textbf{Initialisation (\texttt{gnubg\_init\_rollout()}):} - Allocates
\texttt{rngctxRollout} by copying the main RNG context via
\texttt{CopyRNGContext} - Creates a persistent \texttt{GThreadPool} with
\texttt{g\_get\_num\_processors()} threads - Pool is reused across all
rollout calls to avoid thread creation overhead

\textbf{Thread-local state (two \texttt{GPrivate} keys):} -
\texttt{gnubg\_tls\_key} --- owns \texttt{ThreadLocalData} per thread:
\texttt{aMoves{[}MAX\_MOVES{]}}, \texttt{pnnState{[}3{]}} with
\texttt{savedBase}/\texttt{savedIBase} buffers sized from
\texttt{nnContact} dimensions. Destructor frees all buffers. -
\texttt{gnubg\_rng\_key} --- owns an isolated \texttt{rngcontext} copy
per thread. Copied from \texttt{rngctxRollout} on first use. Destructor
calls \texttt{free\_rngctx()}. This is the critical RNG isolation that
prevents data races.

\textbf{Scatter-gather rollout loop:} 1. Allocate \texttt{nTrials}
result structs with 64-byte cache-line alignment
(\texttt{posix\_memalign}) to prevent false sharing 2. Initialise
\texttt{GMutex}/\texttt{GCond} barrier with
\texttt{tasks\_remaining\ =\ nTrials} 3. Push \texttt{nTrials} tasks to
pool via \texttt{g\_thread\_pool\_push()} 4. Block on
\texttt{g\_cond\_wait()} until all workers signal completion

\textbf{Per-trial worker (\texttt{rollout\_worker\_func}):} -
Retrieves/allocates thread-local \texttt{ThreadLocalData} via
\texttt{TLSGet()} - Retrieves/allocates thread-local \texttt{rngcontext}
via \texttt{gnubg\_rng\_key} - Creates per-task \texttt{perArray} with
seed offset by \texttt{task\_index} for quasi-random dice - Calls
\texttt{BasicCubefulRolloutNoLocking()} with the full 13-argument
signature - Writes result to its designated slot in the scatter array -
Decrements \texttt{tasks\_remaining} and signals \texttt{GCond} when it
reaches zero

\textbf{Scatter-gather merge (main thread after barrier):} - Accumulates
sum and sum-of-squares over all trial results - Computes mean and
standard deviation per output

\paragraph{5.7.3 Rollout globals}\label{rollout-globals}

The following globals are defined in \texttt{stubs.c} with defaults
matching gnubg's standard money-game settings. These are the ``docking
points'' between the engine and the UI layer:

\begin{itemize}
\tightlist
\item
  \texttt{rolloutcontext\ rcRollout} --- default rollout configuration
  (1296 trials, cubeful, variance reduction, quasi-random dice, Mersenne
  Twister RNG)
\item
  \texttt{rngcontext\ *rngctxRollout} --- RNG context for rollouts;
  allocated by \texttt{gnubg\_init\_rollout()} via
  \texttt{CopyRNGContext(rngctxCurrent)}
\item
  \texttt{int\ fAutoCrawford}, \texttt{fAutoSaveRollout},
  \texttt{fShowProgress} --- match/UI state flags
\item
  \texttt{int\ fOutputMWC}, \texttt{fOutputWinPC},
  \texttt{fOutputMatchPC} --- output format flags (declared in
  \texttt{format.h} which is stubbed)
\end{itemize}

\paragraph{5.7.2 gnubg\_init\_rollout()}\label{gnubg_init_rollout}

Must be called after \texttt{EvalInitialise()} and
\texttt{gnubg\_init\_tld()}. Allocates \texttt{rngctxRollout} by copying
the current RNG context established during engine initialisation.
\texttt{rngcontext} is an opaque type defined privately in
\texttt{dice.c}; \texttt{CopyRNGContext} is the only correct way to
create an instance from outside that translation unit.

\paragraph{5.7.3 gnubg\_rollout()}\label{gnubg_rollout}

\begin{Shaded}
\begin{Highlighting}[]
\DataTypeTok{int}\NormalTok{ gnubg\_rollout}\OperatorTok{(}\DataTypeTok{const}\NormalTok{ TanBoard anBoard}\OperatorTok{,}
                  \DataTypeTok{float}\NormalTok{ arOutput}\OperatorTok{[}\NormalTok{NUM\_ROLLOUT\_OUTPUTS}\OperatorTok{],}
                  \DataTypeTok{float}\NormalTok{ arStdDev}\OperatorTok{[}\NormalTok{NUM\_ROLLOUT\_OUTPUTS}\OperatorTok{],}
                  \DataTypeTok{const}\NormalTok{ cubeinfo }\OperatorTok{*}\NormalTok{pci}\OperatorTok{,}\NormalTok{ rolloutcontext }\OperatorTok{*}\NormalTok{prc}\OperatorTok{);}
\end{Highlighting}
\end{Shaded}

Calls \texttt{BasicCubefulRolloutNoLocking} for each of
\texttt{prc-\textgreater{}nTrials} games. Uses
\texttt{QuasiRandomSeed()} (copied from \texttt{rollout.c} where it is
\texttt{static}) to initialise the quasi-random dice permutation array.
Accumulates per-game outputs and computes mean and standard deviation on
completion.

\paragraph{5.7.4 QuasiRandomSeed}\label{quasirandomseed}

\textbf{Note on \texttt{QuasiRandomSeed}:} This function is
\texttt{static} in \texttt{rollout.c} with no header declaration. It
cannot be linked from outside. The function body was copied verbatim
into \texttt{stubs.c}; it uses only \texttt{irandinit}/\texttt{irand}
from \texttt{lib/isaac.c} which is already in the build. The copy is
documented in \texttt{PROVENANCE.md}.

\paragraph{5.7.5 Rollout accuracy}\label{rollout-accuracy}

With \texttt{nTrials=144} (the default minimum) the standard deviation
on win probability is approximately ±0.005. For publication-quality
results use \texttt{nTrials=1296} or higher. The opening position
rollout equity (0.29 at 144 trials) diverges from the 1-ply neural net
estimate (0.077) due to sampling variance; convergence improves with
more trials.

\subsubsection{5.8 android-app.c --- The Android Application
Shell}\label{android-app.c-the-android-application-shell}

\texttt{jni-bridge/src/android-app.c} is the architectural centrepiece
of Phase 8. It replaces \texttt{gnubg.c} (the GTK application shell,
\textasciitilde8000 lines) on Android.

\paragraph{5.8.1 Design Principle}\label{design-principle}

On desktop gnubg, \texttt{gnubg.c} serves two roles: 1. \textbf{Portable
application logic} --- \texttt{InitBoard}, \texttt{NextToken},
\texttt{ParseNumber}, \texttt{find\_skills},
\texttt{GetMatchStateCubeInfo} etc. --- pure C with no GTK dependency 2.
\textbf{GTK application shell} --- \texttt{main()}, GTK initialisation,
window management, signal handlers, readline, Python module, curl

On Android, role 1 is extracted verbatim into \texttt{android-app.c}.
Role 2 is replaced by Android callbacks that post events to the Kotlin
layer. \texttt{progress.c} (GTK rollout progress dialog) and
\texttt{sound.c} are also replaced here.

\paragraph{5.8.2 Android Callback
Architecture}\label{android-callback-architecture}

Six weak-symbol callbacks form the interface between the C engine and
the Android UI layer:

\begin{Shaded}
\begin{Highlighting}[]
\DataTypeTok{void}\NormalTok{ gnubg\_on\_sound\_event}\OperatorTok{(}\NormalTok{gnubgsound gs}\OperatorTok{);}
\DataTypeTok{void}\NormalTok{ gnubg\_on\_setting\_changed}\OperatorTok{(}\DataTypeTok{void} \OperatorTok{*}\NormalTok{pv}\OperatorTok{);}
\DataTypeTok{void}\NormalTok{ gnubg\_on\_board\_changed}\OperatorTok{(}\DataTypeTok{void}\OperatorTok{);}
\DataTypeTok{void}\NormalTok{ gnubg\_on\_navigation\_event}\OperatorTok{(}\DataTypeTok{int}\NormalTok{ ev}\OperatorTok{);}
\DataTypeTok{void}\NormalTok{ gnubg\_on\_filename\_changed}\OperatorTok{(}\DataTypeTok{const} \DataTypeTok{char} \OperatorTok{*}\NormalTok{sz}\OperatorTok{);}
\DataTypeTok{void}\NormalTok{ gnubg\_on\_rollout\_progress}\OperatorTok{(}\DataTypeTok{int}\NormalTok{ iGame}\OperatorTok{,} \DataTypeTok{int}\NormalTok{ nTrials}\OperatorTok{,} \DataTypeTok{float}\NormalTok{ rJsd}\OperatorTok{,} \DataTypeTok{int}\NormalTok{ fStopped}\OperatorTok{);}
\end{Highlighting}
\end{Shaded}

Declared \texttt{\_\_attribute\_\_((weak))} with no-op defaults.
\texttt{native-lib.c} overrides them with implementations that call back
into the JVM.

\paragraph{5.8.3 Stub Policy}\label{stub-policy}

The principle: \textbf{stub only what is genuinely GTK rendering code.}
Everything else is real implementation.

\begin{longtable}[]{@{}
  >{\raggedright\arraybackslash}p{(\linewidth - 4\tabcolsep) * \real{0.3333}}
  >{\raggedright\arraybackslash}p{(\linewidth - 4\tabcolsep) * \real{0.3333}}
  >{\raggedright\arraybackslash}p{(\linewidth - 4\tabcolsep) * \real{0.3333}}@{}}
\toprule\noalign{}
\begin{minipage}[b]{\linewidth}\raggedright
Function
\end{minipage} & \begin{minipage}[b]{\linewidth}\raggedright
Status
\end{minipage} & \begin{minipage}[b]{\linewidth}\raggedright
Future
\end{minipage} \\
\midrule\noalign{}
\endhead
\bottomrule\noalign{}
\endlastfoot
\texttt{playSound} & Routes to \texttt{gnubg\_on\_sound\_event} &
Android SoundPool \\
\texttt{ShowBoard} & Routes to \texttt{gnubg\_on\_board\_changed} &
Android board view \\
\texttt{UpdateSetting} & Routes to \texttt{gnubg\_on\_setting\_changed}
& Android settings binding \\
\texttt{get\_input\_discard} & Returns TRUE & Android confirmation
dialog \\
\texttt{confirmOverwrite} & Returns TRUE & Android file dialog \\
\texttt{GiveAdvice} & Routes to navigation callback & Android tutor
dialog \\
\texttt{GetInputYN} & Returns TRUE & Android yes/no dialog \\
\texttt{CommandFirstGame/Move} & Routes to navigation callback & Android
navigation \\
\texttt{RolloutProgressStart/Progress/End} & Routes to rollout progress
callback & Android progress UI \\
\texttt{HandleCommand} & No-op & Android command layer \\
\texttt{hint\_double/take/move} & No-op & Android tutor hints \\
\texttt{ExtInitParse/ExtStartParse/ExtDestroyParse} & Stubs & External
player protocol \\
\texttt{SetupLanguage} & Returns NULL & Android locale system \\
\end{longtable}

\begin{center}\rule{0.5\linewidth}{0.5pt}\end{center}

\subsection{6. JNI API Reference}\label{jni-api-reference}

\subsubsection{6.1 Board Encoding}\label{board-encoding}

The board is passed as a \texttt{jintArray} of 50 elements encoding
\texttt{TanBoard\ anBoard{[}2{]}{[}25{]}}:

\begin{itemize}
\tightlist
\item
  Elements 0--24: \texttt{anBoard{[}0{]}{[}i{]}} --- opponent's checker
  counts
\item
  Elements 25--49: \texttt{anBoard{[}1{]}{[}i{]}} --- current player's
  checker counts
\end{itemize}

\subsubsection{6.2 Kotlin API (Engine
object)}\label{kotlin-api-engine-object}

\begin{longtable}[]{@{}
  >{\raggedright\arraybackslash}p{(\linewidth - 2\tabcolsep) * \real{0.5000}}
  >{\raggedright\arraybackslash}p{(\linewidth - 2\tabcolsep) * \real{0.5000}}@{}}
\toprule\noalign{}
\begin{minipage}[b]{\linewidth}\raggedright
Method
\end{minipage} & \begin{minipage}[b]{\linewidth}\raggedright
Description
\end{minipage} \\
\midrule\noalign{}
\endhead
\bottomrule\noalign{}
\endlastfoot
\texttt{initialise(weightsPath:\ String):\ Boolean} & Must be called
once before any evaluation. Pass the absolute path to
\texttt{gnubg.weights} extracted to internal storage. Calls
\texttt{EvalInitialise()}, \texttt{SetCubeInfo()},
\texttt{gnubg\_init\_tld()}, and \texttt{gnubg\_init\_rollout()}. \\
\texttt{evaluatePosition(board:\ IntArray):\ FloatArray?} & Returns
\texttt{FloatArray{[}5{]}}:
\texttt{{[}winNormal,\ winGammon,\ winBackgammon,\ loseGammon,\ loseBackgammon{]}}.
Null on engine error. Uses 1-ply cubeless \texttt{evalcontext}. \\
\texttt{findBestMove(board:\ IntArray,\ die0:\ Int,\ die1:\ Int):\ IntArray?}
& Returns \texttt{IntArray{[}8{]}} encoding \texttt{anMove{[}8{]}};
unused slots are -1. Null on error. \\
\texttt{classifyPosition(board:\ IntArray):\ Int} & Returns the
\texttt{positionclass} integer (race, contact, bearoff etc). Returns -1
if not initialised. \\
\texttt{applyMove(board:\ IntArray,\ move:\ IntArray):\ IntArray} &
Applies a move to a board. Returns the resulting 50-element board
encoding. \\
\texttt{cubeDecision(board,\ cubeValue,\ cubeOwner,\ matchTo,\ score0,\ score1,\ crawford):\ IntArray?}
& Returns \texttt{IntArray{[}16{]}}: \texttt{{[}0..6{]}} = if-double
equity floats as \texttt{Int} bits (unpack with
\texttt{Float.fromBits()}), \texttt{{[}7..13{]}} = no-double equity
floats, \texttt{{[}14{]}} = \texttt{CubeDecision} enum value,
\texttt{{[}15{]}} = reserved. See \texttt{Engine.CubeDecision} enum for
all 21 values. \\
\texttt{rollout(board:\ IntArray,\ trials:\ Int\ =\ 144):\ FloatArray?}
& Multi-threaded cubeful rollout via \texttt{GThreadPool}. Returns
\texttt{FloatArray{[}14{]}}: \texttt{{[}0..6{]}} = equity outputs,
\texttt{{[}7..13{]}} = standard deviations. Use \texttt{trials=1296} or
higher for publication-quality results. \\
\texttt{loadSGF(path:\ String):\ Boolean} & Load a match from an SGF
file. Calls \texttt{CommandLoadMatch()} which handles
\texttt{FreeMatch}, \texttt{ClearMatch}, \texttt{RestoreGame},
\texttt{UpdateSettings}. Returns true on success. \\
\texttt{saveSGF(path:\ String):\ Boolean} & Save the current match to an
SGF file. Calls \texttt{CommandSaveMatch()} which handles
\texttt{SaveGame} loop and file I/O. Returns false if no game in
progress. \\
\end{longtable}

\subsubsection{6.3 Thread Safety}\label{thread-safety}

All JNI entry points acquire a static \texttt{pthread\_mutex\_t}
(\texttt{gnubg\_lock}) before calling any engine function and release it
on return. The gnubg evaluation engine uses global state (NNState,
bearoff database handles, weight arrays) that is not safe to access
concurrently.

\subsubsection{6.4 Initialisation \& Data
Files}\label{initialisation-data-files}

The engine requires the following files at runtime, extracted from app
assets to internal storage before calling \texttt{Engine.initialise()}:

\begin{itemize}
\tightlist
\item
  \texttt{gnubg.weights} --- trained neural network weights
  (\textasciitilde6 MB). Primary source of playing strength.
\item
  \texttt{gnubg\_os0.bd} --- small one-sided bearoff database
  (\textasciitilde1.4 MB, 6pt/15ch). Loaded into \texttt{pbc1}.
\item
  \texttt{gnubg\_ts0.bd} --- small two-sided bearoff database
  (\textasciitilde6.8 MB, 6pt/6ch). Loaded into \texttt{pbc2}.
\item
  \texttt{gnubg\_os.bd} --- larger one-sided bearoff database
  (\textasciitilde4.8 MB, 7pt/15ch). Loaded into \texttt{pbcOS}.
  Generated via
  \texttt{makebearoff\ -o\ 7\ -O\ gnubg\_os0.bd\ -f\ gnubg\_os.bd}.
\item
  \texttt{gnubg\_ts.bd} --- two-sided bearoff database
  (\textasciitilde6.6 MB, 6pt/6ch copy). Loaded into \texttt{pbcTS}.
  Copy of \texttt{gnubg\_ts0.bd}.
\end{itemize}

All four databases are deployed and confirmed loading on device
(\texttt{pbc1}, \texttt{pbc2}, \texttt{pbcOS}, \texttt{pbcTS} all
non-null). The \texttt{AC\_PKGDATADIR} path is compiled in as
\texttt{/data/data/com.clavierhaus.gnubg/files} and must match the
actual extraction location.

\begin{center}\rule{0.5\linewidth}{0.5pt}\end{center}

\subsection{7. Build Instructions}\label{build-instructions}

\subsubsection{7.1 Prerequisites (Fedora
44)}\label{prerequisites-fedora-44}

All prerequisites are available via dnf. No pip, npm, or external
package managers are required:

\begin{Shaded}
\begin{Highlighting}[]
\FunctionTok{sudo}\NormalTok{ dnf install meson ninja{-}build cmake glib2{-}devel}
\end{Highlighting}
\end{Shaded}

NDK location:
\texttt{/home/erweitert/gnubg-android/android-sdk/ndk/27.0.11718014/}

\subsubsection{7.2 Build GLib for Android (one-time, \textasciitilde3--5
minutes)}\label{build-glib-for-android-one-time-35-minutes}

\begin{Shaded}
\begin{Highlighting}[]
\BuiltInTok{cd}\NormalTok{ /home/erweitert/gnubg{-}android}
\ExtensionTok{./build\_glib\_android.sh} \DecValTok{2}\OperatorTok{\textgreater{}\&}\DecValTok{1} \KeywordTok{|} \FunctionTok{tee}\NormalTok{ glib{-}build.log}
\end{Highlighting}
\end{Shaded}

On success: \texttt{jni-bridge/external/glib/lib/libglib-2.0.so} is
created (ELF 64-bit ARM aarch64, Android 28).

\subsubsection{7.3 Build
libgnubg-engine.so}\label{build-libgnubg-engine.so}

\begin{Shaded}
\begin{Highlighting}[]
\BuiltInTok{cd}\NormalTok{ /home/erweitert/gnubg{-}android/jni{-}bridge}
\FunctionTok{rm} \AttributeTok{{-}rf}\NormalTok{ build }\KeywordTok{\&\&} \FunctionTok{mkdir}\NormalTok{ build }\KeywordTok{\&\&} \BuiltInTok{cd}\NormalTok{ build}
\FunctionTok{cmake}\NormalTok{ .. }\DataTypeTok{\textbackslash{}}
  \AttributeTok{{-}DCMAKE\_TOOLCHAIN\_FILE}\OperatorTok{=}\NormalTok{/home/erweitert/gnubg{-}android/android{-}sdk/ndk/27.0.11718014/build/cmake/android.toolchain.cmake }\DataTypeTok{\textbackslash{}}
  \AttributeTok{{-}DANDROID\_ABI}\OperatorTok{=}\NormalTok{arm64{-}v8a }\DataTypeTok{\textbackslash{}}
  \AttributeTok{{-}DANDROID\_PLATFORM}\OperatorTok{=}\NormalTok{android{-}28 }\DataTypeTok{\textbackslash{}}
  \AttributeTok{{-}DCMAKE\_BUILD\_TYPE}\OperatorTok{=}\NormalTok{Release}
\FunctionTok{cmake} \AttributeTok{{-}{-}build}\NormalTok{ . }\DecValTok{2}\OperatorTok{\textgreater{}\&}\DecValTok{1} \KeywordTok{|} \FunctionTok{tee}\NormalTok{ /home/erweitert/gnubg{-}android/jni{-}build.log}
\end{Highlighting}
\end{Shaded}

\subsubsection{7.4 Verification}\label{verification}

\begin{Shaded}
\begin{Highlighting}[]
\FunctionTok{file}\NormalTok{ jni{-}bridge/build/libgnubg{-}engine.so}
\CommentTok{\# Expected: ELF 64{-}bit LSB shared object, ARM aarch64, for Android 28}

\FunctionTok{nm}\NormalTok{ jni{-}bridge/build/libgnubg{-}engine.so }\KeywordTok{|} \FunctionTok{grep} \StringTok{"T Java\_"}
\CommentTok{\# Expected (9 JNI entry points):}
\CommentTok{\#   T Java\_com\_clavierhaus\_gnubg\_Engine\_applyMove}
\CommentTok{\#   T Java\_com\_clavierhaus\_gnubg\_Engine\_classifyPosition}
\CommentTok{\#   T Java\_com\_clavierhaus\_gnubg\_Engine\_cubeDecision}
\CommentTok{\#   T Java\_com\_clavierhaus\_gnubg\_Engine\_evaluatePosition}
\CommentTok{\#   T Java\_com\_clavierhaus\_gnubg\_Engine\_findBestMove}
\CommentTok{\#   T Java\_com\_clavierhaus\_gnubg\_Engine\_initialise}
\CommentTok{\#   T Java\_com\_clavierhaus\_gnubg\_Engine\_loadSGF}
\CommentTok{\#   T Java\_com\_clavierhaus\_gnubg\_Engine\_rollout}
\CommentTok{\#   T Java\_com\_clavierhaus\_gnubg\_Engine\_saveSGF}

\FunctionTok{nm}\NormalTok{ jni{-}bridge/build/libgnubg{-}engine.so }\KeywordTok{|} \FunctionTok{grep} \StringTok{"NeuralNetEvaluateSSE"}
\CommentTok{\# Expected: T NeuralNetEvaluateSSE  (confirms NEON path active)}

\FunctionTok{nm}\NormalTok{ jni{-}bridge/build/libgnubg{-}engine.so }\KeywordTok{|} \FunctionTok{grep} \StringTok{"gnubg\_rollout\textbackslash{}|BasicCubefulRollout"}
\CommentTok{\# Expected: T gnubg\_rollout, T BasicCubefulRolloutNoLocking}
\end{Highlighting}
\end{Shaded}

\subsubsection{7.5 Android adb Test
Harness}\label{android-adb-test-harness}

The test harness in \texttt{test-harness/} builds as a standalone
Android executable and exercises the engine end-to-end on a connected
device.

\begin{Shaded}
\begin{Highlighting}[]
\CommentTok{\# Build}
\BuiltInTok{cd}\NormalTok{ /home/erweitert/gnubg{-}android/test{-}harness}
\FunctionTok{rm} \AttributeTok{{-}rf}\NormalTok{ build }\KeywordTok{\&\&} \FunctionTok{mkdir}\NormalTok{ build }\KeywordTok{\&\&} \BuiltInTok{cd}\NormalTok{ build}
\FunctionTok{cmake}\NormalTok{ .. }\DataTypeTok{\textbackslash{}}
  \AttributeTok{{-}DCMAKE\_TOOLCHAIN\_FILE}\OperatorTok{=}\NormalTok{/home/erweitert/gnubg{-}android/android{-}sdk/ndk/27.0.11718014/build/cmake/android.toolchain.cmake }\DataTypeTok{\textbackslash{}}
  \AttributeTok{{-}DANDROID\_ABI}\OperatorTok{=}\NormalTok{arm64{-}v8a }\DataTypeTok{\textbackslash{}}
  \AttributeTok{{-}DANDROID\_PLATFORM}\OperatorTok{=}\NormalTok{android{-}28 }\DataTypeTok{\textbackslash{}}
  \AttributeTok{{-}DCMAKE\_BUILD\_TYPE}\OperatorTok{=}\NormalTok{Release}
\FunctionTok{cmake} \AttributeTok{{-}{-}build}\NormalTok{ .}

\CommentTok{\# Push to device}
\ExtensionTok{adb}\NormalTok{ push test\_runner /data/local/tmp/gnubg{-}test\_runner}
\ExtensionTok{adb}\NormalTok{ push /path/to/gnubg.weights   /data/local/tmp/gnubg.weights}
\ExtensionTok{adb}\NormalTok{ push /path/to/gnubg\_os0.bd   /data/local/tmp/gnubg\_os0.bd}
\ExtensionTok{adb}\NormalTok{ push /path/to/gnubg\_ts0.bd   /data/local/tmp/gnubg\_ts0.bd}
\ExtensionTok{adb}\NormalTok{ push jni{-}bridge/external/glib/lib/libglib{-}2.0.so /data/local/tmp/}
\ExtensionTok{adb}\NormalTok{ push jni{-}bridge/external/glib/lib/libintl.so     /data/local/tmp/}
\ExtensionTok{adb}\NormalTok{ shell chmod +x /data/local/tmp/gnubg{-}test\_runner}

\CommentTok{\# Run}
\ExtensionTok{adb}\NormalTok{ shell }\StringTok{"LD\_LIBRARY\_PATH=/data/local/tmp }\DataTypeTok{\textbackslash{}}
\StringTok{    /data/local/tmp/gnubg{-}test\_runner /data/local/tmp/gnubg.weights"}
\end{Highlighting}
\end{Shaded}

Expected output:

\begin{verbatim}
=== GNU Backgammon Android Test Harness ===
[1] Initialising engine...
    pbc1=0x... pbc2=0x... pbcOS=0x... pbcTS=0x...  (all non-null)
[2] ClassifyPosition = 10 (expected 10 = CLASS_CONTACT)
[3] EvaluatePosition 1-ply cubeless...
    Win prob:        0.5269
    Win gammon:      0.1478
    ...
    Cubeless equity: 0.0768
[4] FindBestMove for 3-1...
    Move: 8/5 6/5 (expected 8/5 6/5)
[5] Cube decision (money game, centred cube)...
    Cube decision: 2 (NODOUBLE_TAKE)
    No double equity:   0.4260
    Double/take equity: 0.4260
    Double/pass equity: 0.2950
[6] Rollout (144 trials, cubeful variance reduction)...
    gnubg_rollout returned 0
    Win prob:  ~0.59  stddev: ~0.005
    Equity:    ~0.29  (converges toward 1-ply with more trials)
=== All tests passed ===
\end{verbatim}

\begin{center}\rule{0.5\linewidth}{0.5pt}\end{center}

\subsection{8. Open Items \& Roadmap}\label{open-items-roadmap}

\subsubsection{8.1 Features}\label{features}

The engine is feature-complete as of Phase 9.

\begin{itemize}
\tightlist
\item
  \textbf{Rollout:} {[}DONE{]} Multi-threaded via
  \texttt{GThreadPool}/\texttt{GPrivate}. \texttt{Engine.rollout()}. See
  §5.7.
\item
  \textbf{Cube decisions:} {[}DONE{]} \texttt{Engine.cubeDecision()}.
  Full equity matrix + \texttt{CubeDecision} enum (21 values).
\item
  \textbf{SGF import/export:} {[}DONE{]} \texttt{Engine.loadSGF()} /
  \texttt{Engine.saveSGF()}. Backed by
  \texttt{CommandLoadMatch}/\texttt{CommandSaveMatch} in \texttt{sgf.c}.
\item
  \textbf{Game management:} {[}DONE{]} \texttt{play.c} compiled. Full
  match state machine.
\item
  \textbf{Analysis:} {[}DONE{]} \texttt{analysis.c} compiled.
  \texttt{find\_skills}, \texttt{AnalyzeMove}, \texttt{AnalyzeGame},
  \texttt{AnalyzeMatch} available.
\item
  \textbf{Multi-threading:} {[}DONE{]} \texttt{USE\_MULTITHREAD}
  enabled. \texttt{GThreadPool} + \texttt{GPrivate} TLS + RNG isolation.
  See §5.7.
\item
  \textbf{SQLite:} Not planned --- use Android's Room/SQLiteOpenHelper
  for persistence.
\end{itemize}

\subsubsection{8.2 Immediate Next Steps}\label{immediate-next-steps}

\begin{itemize}
\tightlist
\item
  \textbf{Android Studio project:} Create the Android app project, add
  \texttt{libgnubg-engine.so} and GLib \texttt{.so} files as prebuilt
  native libraries, add \texttt{Engine.kt} to the source set.
\item
  \textbf{Asset packaging:} Add \texttt{gnubg.weights} and all four
  bearoff databases to the app's assets. Implement extraction to
  \texttt{context.filesDir} on first launch.
\item
  \textbf{Callback wiring:} Implement non-weak versions of
  \texttt{gnubg\_on\_*} callbacks in \texttt{native-lib.c} to route
  engine events into the Kotlin UI layer.
\end{itemize}

\subsubsection{8.3 Performance}\label{performance}

\begin{itemize}
\tightlist
\item
  \textbf{ARM NEON SIMD:} Enabled. \texttt{NeuralNetEvaluateSSE}
  confirmed exported and active. \texttt{CheckNEON()} always returns 1
  on aarch64.
\item
  \textbf{Multi-threading:} {[}DONE{]} \texttt{USE\_MULTITHREAD}
  enabled. Rollout tasks distributed across all CPU cores via
  \texttt{GThreadPool}. Per-thread \texttt{rngcontext} isolation via
  \texttt{GPrivate} guarantees correctness. See §5.7.
\end{itemize}

\subsubsection{8.4 Known Warnings}\label{known-warnings}

\begin{itemize}
\tightlist
\item
  \textbf{\texttt{multithread.c}:} 6 instances of incompatible pointer
  types passing \texttt{pthread\_mutex\_t\ *} to \texttt{Mutex\ *}.
  \texttt{USE\_MULTITHREAD} is now enabled but these coercions are in
  paths that are compiled but not reached on Android; harmless in
  practice.
\item
  \textbf{CMake deprecation warnings} from the NDK toolchain file
  regarding
  \texttt{cmake\_minimum\_required\ VERSION\ \textless{}\ 3.10}. In the
  NDK's own toolchain file; cannot be fixed from the project.
\item
  \textbf{Rollout equity variance at low trial counts:} With
  \texttt{nTrials=144} stddev on win prob is \textasciitilde0.005. Use
  1296+ trials for reliable equity estimates.
\end{itemize}

\begin{center}\rule{0.5\linewidth}{0.5pt}\end{center}

\subsection{9. Key File Reference}\label{key-file-reference}

\subsubsection{build\_glib\_android.sh}\label{build_glib_android.sh}

Orchestrates the complete GLib cross-build: applies Fedora SRPM patches,
patches \texttt{meson.build} for Android iconv-in-libc, runs
\texttt{meson\ setup} with \texttt{android-arm64.cross}, builds with
ninja, installs to \texttt{jni-bridge/external/glib/}. Safe to re-run.

\subsubsection{android-arm64.cross}\label{android-arm64.cross}

Meson cross-file for the GLib build. Specifies NDK 27 toolchain binaries
(\texttt{aarch64-linux-android28-clang}), target machine
(android/aarch64/little-endian), compile flags
(\texttt{-DANDROID\ -fPIC\ -Os}), and platform size properties. API
level 28 is encoded in the compiler binary name, not as a \texttt{-D}
flag.

\subsubsection{jni-bridge/CMakeLists.txt}\label{jni-bridgecmakelists.txt}

NDK CMake build file for \texttt{libgnubg-engine.so}. Enumerates all
compiled source files, sets include paths (\texttt{jni-bridge/} first
for shadow \texttt{config.h} interception, then \texttt{engine-core/},
\texttt{engine-core/lib/}, GLib headers), defines
\texttt{AC\_DATADIR}/\texttt{AC\_PKGDATADIR}/\texttt{AC\_DOCDIR} for
Android, and links against \texttt{libglib-2.0.so},
\texttt{libgobject-2.0.so}, \texttt{libintl.so}, \texttt{liblog}, and
\texttt{libm}.

\subsubsection{jni-bridge/src/stubs.c}\label{jni-bridgesrcstubs.c}

Provides remaining GTK/UI stub storage, TLD infrastructure (now via
\texttt{GPrivate}), multi-threaded rollout orchestration
(\texttt{gnubg\_rollout()} with \texttt{GThreadPool}), and rollout
globals. Key functions: \texttt{gnubg\_init\_tld()},
\texttt{gnubg\_init\_rollout()}, \texttt{gnubg\_rollout()},
\texttt{QuasiRandomSeed()}, \texttt{TLSGet()}. \textbf{Note:}
\texttt{test-harness/src/stubs.c} is a copy of this file and must be
kept in sync.

\subsubsection{jni-bridge/src/android-app.c}\label{jni-bridgesrcandroid-app.c}

The Android application shell replacing \texttt{gnubg.c}. Contains
functions extracted verbatim from \texttt{gnubg.c} (see §5.8) plus
Android callbacks for sound, board, settings, navigation, and rollout
progress. All globals previously in \texttt{gnubg.c} are defined here.

\subsubsection{engine-core/config.h}\label{engine-coreconfig.h}

Host-generated autotools \texttt{config.h}, regenerated by
\texttt{grep\ -v} to strip Android-incompatible flags.
\textbf{Important:} if \texttt{engine-core/} is refreshed from upstream,
this file must be regenerated and re-patched using the \texttt{grep\ -v}
pipeline documented in §5.3.

\subsubsection{engine-core/lib/neuralnetsse.c}\label{engine-corelibneuralnetsse.c}

Modified from upstream: \texttt{NeuralNetEvaluateSSE} function has its
VLA replaced with \texttt{posix\_memalign} allocation for correct
16-byte alignment on aarch64. See §5.6 and \texttt{PROVENANCE.md} for
full documentation of this change.

\begin{center}\rule{0.5\linewidth}{0.5pt}\end{center}

\subsection{10. Android UI --- IDE Decision, Design Philosophy \& Way
Forward}\label{android-ui-ide-decision-design-philosophy}

\subsubsection{10.1 Milestone: Engine
Complete}\label{milestone-engine-complete}

As of Phase 9, \texttt{libgnubg-engine.so} is feature-complete. All 9
JNI entry points are implemented and verified on a Pixel 8 Pro (aarch64,
Android 16). The engine milestone closes the native layer of the project.
Version v0.9 marks the transition from engine development to Android UI
development, targeting a 1.0 release.

\subsubsection{10.2 IDE Decision: No IDE --- Pure Terminal
Workflow}\label{ide-decision-no-ide}

After evaluating all options, the decision is to use a \textbf{pure
terminal workflow} with no Android Studio or other IDE involved at any
stage. The rationale:

\begin{itemize}
\tightlist
\item
  Android Studio requires a GUI session, significant RAM overhead, and
  account/sign-in infrastructure that adds friction without benefit for
  this project.
\item
  The entire build is driven by \texttt{gradle assembleDebug} from the
  terminal --- the same Gradle invocation Android Studio uses internally.
\item
  All source files are plain Kotlin and are edited in any text editor.
  The project has no XML layouts (see §10.3).
\item
  The keyboard-operator workflow (human pastes commands, observes
  output) is better served by explicit, auditable shell commands than by
  IDE wizards.
\end{itemize}

The project structure is a bare Gradle project (\texttt{gnubg-app/})
created entirely from the terminal. No \texttt{.idea/} directory, no
Gradle wrapper magic, no IDE-generated boilerplate beyond what is
explicitly documented here.

\subsubsection{10.3 UI Toolkit: Jetpack
Compose}\label{ui-toolkit-jetpack-compose}

The Android UI is built exclusively with \textbf{Jetpack Compose} ---
Android's modern declarative UI toolkit. XML layouts are explicitly
excluded. Rationale:

\begin{itemize}
\tightlist
\item
  Compose maps directly to the terminal/text-editor workflow: every UI
  element is a Kotlin \texttt{@Composable} function. There are no
  parallel XML resource files to maintain.
\item
  Material You dynamic colour (\texttt{dynamicColorScheme()}) is
  first-class in Compose. All colour tokens defined in §10.5 are
  available as \texttt{MaterialTheme.colorScheme.*} with no additional
  plumbing.
\item
  The board is drawn via Compose's \texttt{Canvas} composable --- the
  same drawing model as Cairo/SVG, accepting the boarddim.h geometry
  directly.
\item
  Animation (checker movement, dice roll) is \texttt{ObjectAnimator}
  over Compose state --- no separate animation framework needed.
\end{itemize}

\subsubsection{10.4 Build Configuration}\label{build-configuration}

\begin{longtable}[]{@{}ll@{}}
\toprule\noalign{}
Parameter & Value \\
\midrule\noalign{}
\endhead
\bottomrule\noalign{}
\endlastfoot
\texttt{minSdk} & 31 (Android 12 --- dynamic colour guaranteed on all supported devices) \\
\texttt{targetSdk} & 35 (Android 15 --- current Play Store requirement) \\
\texttt{compileSdk} & 35 \\
Kotlin & 2.0.21 \\
Compose BOM & 2024.09.03 (ships Compose 1.7) \\
AGP & 8.5.2 \\
Gradle wrapper & 8.7 \\
\end{longtable}

\texttt{minSdk = 31} means \texttt{dynamicColorScheme()} is available on
every supported device with no API version guards. Play Store reach at
Android 12+ is approximately 75\% of active devices --- acceptable for a
serious backgammon application not targeting the mass casual market.

\subsubsection{10.5 Design Philosophy}\label{design-philosophy}

\paragraph{Orientation}\label{orientation}

Landscape only. No portrait mode is available or planned.
\texttt{android:screenOrientation="landscape"} is set in
\texttt{AndroidManifest.xml}. This fixes the aspect ratio and eliminates
all layout reflow complexity.

\paragraph{Board Geometry}\label{board-geometry}

The board geometry is derived directly from \texttt{boarddim.h} in the
gnubg source. The canonical coordinate system is a 108×82 unit grid
(aspect ratio 1.317:1). All element positions --- point triangles, bar,
bearoff trays, checker stacks, cube, hinges --- are computed from the
\texttt{POINT\_X()}, \texttt{TOP\_POINT\_Y}, \texttt{BOT\_POINT\_Y},
\texttt{CHEQUER\_HEIGHT}, and \texttt{CHEQUER\_WIDTH} macros already
present in the source. This is a one-time translation; the geometry is
gnubg's own, faithfully reproduced.

\paragraph{Material You Colour
System}\label{material-you-colour-system}

The board colour system is built entirely on Android's Material You
dynamic colour API (\texttt{DynamicColors}, \texttt{dynamicColorScheme}).
The device wallpaper seed generates the full tonal palette
automatically. Three user-selectable seed overrides are offered:

\begin{longtable}[]{@{}lll@{}}
\toprule\noalign{}
Name & Seed & Character \\
\midrule\noalign{}
\endhead
\bottomrule\noalign{}
\endlastfoot
Dynamic & Device wallpaper & Default --- fully personalised \\
Classic & \texttt{\#5D4037} & Warm brown → amber/ochre tones \\
Ocean & \texttt{\#1565C0} & Deep blue (BG Galaxy territory) \\
\end{longtable}

Colour role assignments for board elements:

\begin{longtable}[]{@{}ll@{}}
\toprule\noalign{}
Element & Material You Role \\
\midrule\noalign{}
\endhead
\bottomrule\noalign{}
\endlastfoot
Board field & \texttt{tertiaryContainer} \\
Triangle A (dominant) & \texttt{tertiary} \\
Triangle B (recessive) & \texttt{tertiaryContainer} at 60\% \\
Bar & \texttt{surface} darkened \\
Border/frame & \texttt{secondaryContainer} \\
Point numbers & \texttt{onSecondaryContainer} \\
Dark checker fill & \texttt{primary} at P-20 \\
Light checker fill & \texttt{primary} at P-90 \\
Checker highlight & white at 25--30\% opacity, offset top-left \\
Cube fill & \texttt{secondaryContainer} \\
Cube number & \texttt{onSecondaryContainer} \\
Die fill & \texttt{secondary} \\
Die pips & \texttt{onSecondary} \\
\end{longtable}

\textbf{Rationale for tertiary board / primary checkers:} Using
\texttt{tertiary} for the board field and \texttt{primary} for the
checkers guarantees perceptual distinctness across all wallpaper seeds,
because Material You's algorithm explicitly assigns Primary and Tertiary
to different hue regions. Both checker colours (P-20 and P-90) are
within the same Primary hue family, giving a coherent monochromatic
checker aesthetic while remaining legible against the Tertiary board.

\paragraph{Depth Treatment}\label{depth-treatment}

Checkers, cube, and dice share a consistent, minimal depth language:

\begin{enumerate}
\def\labelenumi{\arabic{enumi}.}
\tightlist
\item
  \textbf{Base shape} --- circle (checkers) or rounded-rect (cube/dice),
  solid fill.
\item
  \textbf{Rim} --- 1.5px stroke, one tonal step darker than fill.
\item
  \textbf{Highlight ellipse} --- white at 25--30\% opacity, positioned
  at \texttt{(cx\ -\ 0.28r,\ cy\ -\ 0.28r)}, sized at
  \texttt{0.76r\ ×\ 0.38r}. Simulates a diffuse overhead-left light
  source.
\end{enumerate}

No drop shadows, no gradients, no bevels. The result reads as material
without crossing into skeuomorphic excess.

\paragraph{Point Numbers}\label{point-numbers}

Point numbers are placed \textbf{outside} the board border, in the
frame area above and below the playing surface. They do not appear on
the playing surface. This costs less than 5\% of screen width and
eliminates distraction during play. Players who do not use point numbers
are unaffected.

\paragraph{UI Layer Architecture}\label{ui-layer-architecture}

The Compose UI separates into two layers:

\begin{itemize}
\tightlist
\item
  \textbf{Static board layer} --- a single \texttt{Canvas} composable
  drawing the board field, triangles, bar, bearoff trays, border, and
  point numbers. Drawn once; does not change during a game.
\item
  \textbf{Dynamic piece layer} --- checkers, dice, cube, legal-move
  highlights, and last-move trail. Positioned programmatically from
  \texttt{ChequerPosition()} and \texttt{CubePosition()} in
  \texttt{boardpos.c} via JNI. Checker movement is an
  \texttt{ObjectAnimator} over Compose state.
\end{itemize}

This separation means animations touch only the dynamic layer; the board
surface is never redrawn during play.

\paragraph{Functionality Philosophy}\label{functionality-philosophy}

The UI presents a clean, modern surface while exposing 100\% of gnubg's
analytical depth. Two access modes:

\begin{itemize}
\tightlist
\item
  \textbf{Standard mode} --- board, dice, cube controls, score display.
  Clean and uncluttered.
\item
  \textbf{Expert/Advanced mode} --- analysis panel showing equity, error
  classification, rollout button, hint mode, tutor mode, SGF
  import/export. Everything the desktop version offers, surfaced
  progressively.
\end{itemize}

No gnubg functionality is removed or stubbed at the UI level. The
distinction is presentation only.

\subsubsection{10.6 Immediate Next Steps}\label{immediate-next-steps-v09}

\begin{enumerate}
\def\labelenumi{\arabic{enumi}.}
\tightlist
\item
  Create the bare Gradle project structure (\texttt{gnubg-app/}) from
  the terminal --- \texttt{settings.gradle.kts},
  \texttt{build.gradle.kts}, \texttt{gradle/libs.versions.toml},
  \texttt{AndroidManifest.xml}, \texttt{MainActivity.kt},
  \texttt{ui/theme/Theme.kt}.
\item
  Verify \texttt{gradle assembleDebug} produces a clean APK with an
  empty Compose screen.
\item
  Implement \texttt{Board.kt} --- the static board \texttt{Canvas}
  composable from boarddim.h geometry.
\item
  Implement \texttt{BoardState.kt} and \texttt{EngineViewModel.kt} ---
  the bridge between the JNI engine and Compose state.
\item
  Wire checkers, dice, cube as the dynamic piece layer.
\item
  Add the analysis panel as the Expert mode overlay.
\end{enumerate}

\begin{center}\rule{0.5\linewidth}{0.5pt}\end{center}

\emph{GNU Backgammon Android Port --- MASTER v0.9 --- clavierhaus.at ---
June 2026}

\end{document}
